% Pregunta 983 (2026-marzo, n.º 19): modelo q con costes de ajuste y depreciación nula.
% Aumento permanente de la renta: el locus qdot=0 sube, q salta a la nueva senda de silla
% y K crece hasta K1 mientras q vuelve a 1 (la inversión vuelve a 0).
% Linealización a escala (unidades del dibujo): Kdot = (q-1), qdot = 0,5(q-1) + 0,5(K-K*)
%   qdot=0: pendiente -1; senda de silla: autovalor -0,5, pendiente -0,5 (más plana).
%   E0 = (2, 1)   E1 = (4, 1)   salto: q pasa de 1 a 2 en K0 (punto A)
% Autor: Víctor Gutiérrez Marcos
\documentclass[tikz,border=6pt]{standalone}
\input{tikz-tcee}
\begin{document}
\begin{tikzpicture}[tcee,x=1.1cm,y=1.1cm]
  \ejes{$K$}{$q$}{6.2}{4.8}
  \node[font=\small\bfseries,text=tceemorado] at (3.1,5.2) {Aumento permanente de la renta ($\delta=0$)};
  % Kdot = 0: q = 1
  \draw[curva2] (0,2) -- (6,2) node[etiqueta,above left,text=tceeazul] {$\dot K=0$};
  \node[left] at (0,2) {$1$};
  % Locus qdot = 0 inicial (q = 2 - (K-2) en el dibujo) y desplazado
  \draw[curva1] (0.9,3.1) -- (3.7,0.3) node[etiqueta,below,text=tceemorado] {$\dot q=0$};
  \draw[curvanueva,desplazada] (2.4,3.6) -- (5.7,0.3) node[etiqueta,right,text=tceerojo] {$\dot q=0'$};
  % Sendas de silla: inicial (y = 2 - 0,5(K-2)) y nueva (y = 2 - 0,5(K-4))
  \draw[curva3,dashed] (0.4,2.8) -- (3.6,1.2) node[etiqueta,below right,text=tceeverde] {$SS$};
  \draw[curvanueva] (0.6,3.7) -- (5.8,1.1) node[etiqueta,right,text=tceerojo] {$SS'$};
  % Proyecciones
  \draw[guia] (2,2) -- (2,0) node[below] {$K_0$};
  \draw[guia] (4,2) -- (4,0) node[below] {$K_1$};
  % Salto de q y recorrido por la nueva senda
  \draw[->,very thick,black] (2,2.06) -- (2,2.9);
  \draw[->,very thick,black] (2.4,2.66) -- (3.0,2.36);
  \draw[->,very thick,black] (3.15,2.285) -- (3.6,2.06);
  % Puntos
  \node[punto] at (2,2) {};
  \node[etiqueta,below left] at (2,2) {$E_0$};
  \node[punto] at (2,3) {};
  \node[etiqueta,above right] at (2,3) {$A$};
  \node[punto] at (4,2) {};
  \node[etiqueta,below left] at (4,2) {$E_1$};
\end{tikzpicture}
\end{document}
